\appendix

\section{Proof Details}
\label{sec:proof-details}

\subsection{Boundary cases and the architecture count}
\label{app:step-001}

\begin{proposition}[Exact empty-class closure]
\label{prop:step-001-empty}
Under Assumption~\ref{assump:source-regime} and the convention
$\operatorname{dc}(\varnothing)=0$, if $\mathcal H=\varnothing$, then
\[
\operatorname{dc}(\mathcal H)=0\leq7TSd.
\]
\end{proposition}

\begin{proof}
The convention gives
$\operatorname{dc}(\mathcal H)=\operatorname{dc}(\varnothing)=0$.
Assumption~\ref{assump:source-regime} gives $T\geq1$, $d\geq0$, and a
nonempty sum $S=\sum_{i=1}^L n_i n_{i-1}$ of positive integer products.
Hence $7TSd\geq0$.
\end{proof}

\begin{lemma}[Zero-dimensional tie-resolved closure]
\label{lem:step-001-zero}
Under Assumptions~\ref{assump:source-regime} and
\ref{assump:tie-resolved-confident-map} and the conventions
$\mathbb R^0=\{0\}$, $\langle0,0\rangle=0$, and
$\operatorname{sgn}_{\tau}(0)=\tau$, if $\mathcal H\neq\varnothing$ and
$d=0$, then every $h\in\mathcal H$ satisfies $h(x)=\tau$ for all
$x\in\mathcal X$, and
\[
\operatorname{dc}(\mathcal H)=0=7TSd.
\]
\end{lemma}

\begin{proof}
Since $\mathbb R^0$ is the singleton $\{0\}$, there is exactly one map from
$\mathcal X$ to $\mathbb R^0$, namely $\phi_0(x)=0$ for every $x$.
Consequently every probability law on these maps, including $\mathcal P$ in
Assumption~\ref{assump:tie-resolved-confident-map}, assigns probability one
to $\phi_0$.  There is also exactly one separator in $\mathbb R^0$, namely
$w=0$.

Fix $h\in\mathcal H$.  Because $n\geq1$, the domain
$\mathcal X=\{-1,+1\}^n$ is nonempty, so the universal distribution
quantifier in Assumption~\ref{assump:tie-resolved-confident-map} may be
instantiated at a point mass.  For this fixed target define
\[
E_h:=\left\{\phi:\ \exists w\in\mathbb R^0\ \forall x\in\mathcal X,
\ \operatorname{sgn}_{\tau}(\langle w,\phi(x)\rangle)=h(x)\right\}.
\]
For the sole possible map and separator,
\[
\begin{aligned}
\phi_0\in E_h
&\iff
\forall x\in\mathcal X,
\ \operatorname{sgn}_{\tau}(\langle0,0\rangle)=h(x)\\
&\iff
\forall x\in\mathcal X,
\ h(x)=\tau.
\end{aligned}
\]
Thus $E_h$ is deterministic: its $\mathcal P$-probability is either zero or
one.  Assumption~\ref{assump:tie-resolved-confident-map} gives
$\mathcal P(E_h)\geq1/2$, so $\mathcal P(E_h)=1$ and $h(x)=\tau$ for every
$x$.  Since $h$ was arbitrary, every member of $\mathcal H$ is the constant
$\tau$ classifier.

The same map $\phi_0$, together with $u_h=0$ for every target, satisfies
\[
\operatorname{sgn}_{\tau}(\langle u_h,\phi_0(x)\rangle)
=\operatorname{sgn}_{\tau}(0)
=\tau
=h(x)
\qquad
\text{for all }h\in\mathcal H,
\ x\in\mathcal X.
\]
Dimension zero is therefore feasible in the definition of
$\operatorname{dc}(\mathcal H)$.  Since the minimum is over nonnegative
integer dimensions, $\operatorname{dc}(\mathcal H)=0$.  Finally $d=0$
gives $7TSd=0$.
\end{proof}

\begin{proposition}[First-layer structural bound]
\label{prop:step-001-architecture}
Under Assumption~\ref{assump:source-regime}, if
$\mathcal H\neq\varnothing$ and $d\neq0$, then
\[
d\geq1,
\qquad
S\geq n\geq1,
\qquad
T\geq1,
\qquad
S\geq1.
\]
Moreover, if $L=1$, then $S=n$.
\end{proposition}

\begin{proof}
Assumption~\ref{assump:source-regime} states that
$d\in\mathbb Z_{\geq0}$, so $d\neq0$ implies $d\geq1$.  It also gives
$n,L,T\in\mathbb Z_{\geq1}$, hence $n,T\geq1$.

For every allowed depth, $n_1\geq1$.  When $L\geq2$, $n_1$ is a positive
hidden width, while when $L=1$, it is the output width $n_1=n_L=1$.  Thus
the first summand in the parameter count satisfies
\[
n_1n_0=n_1n\geq n.
\]
Every summand in $S=\sum_{i=1}^L n_i n_{i-1}$ is positive, so
\[
S\geq n_1n_0\geq n\geq1.
\]
If $L=1$, the sum consists only of this first term and $n_1=1$, whence
\[
S=n_1n_0=1\cdot n=n.
\]
This includes the depth-one boundary without positing a hidden layer.
\end{proof}

\begin{proof}[Boundary reduction]
The case split is exhaustive.  If $\mathcal H=\varnothing$, Proposition
\ref{prop:step-001-empty} proves the desired inequality exactly.  Otherwise
$\mathcal H\neq\varnothing$.  If $d=0$, Lemma
\ref{lem:step-001-zero} proves
$\operatorname{dc}(\mathcal H)=0=7TSd$ with the fixed tie convention.  If
$d\neq0$, Proposition~\ref{prop:step-001-architecture} gives
\[
\mathcal H\neq\varnothing,
\qquad
d\geq1,
\qquad
S\geq n\geq1,
\qquad
T,S\geq1,
\]
including $L=1$.  These are the structural controls used in the remaining
proof.
\end{proof}
